\section{The Triple Foundation of the Social Order}

Rene Guenon assigned a triple foundation to the social order and political power:

\begin{enumerate}


\item \textbf{Metaphysical}: \emph{The doctrine of the two powers}.


\item \textbf{Cosmological}: \emph{The doctrine of cycles}.


\item \textbf{Sociological}: \emph{The doctrine of castes}\footnote{\url{https://gornahoor.net/?p=713}}.

\end{enumerate}


The doctrine of the two powers refers to the relationship between \textbf{Spiritual Authority} and \textbf{Temporal Power}. Political power must be subservient to spiritual Wisdom.

The doctrine of cycles refers to the four ages, the decline associated with the transitions between ages, and the recovery of the \textbf{Primordial Tradition} and \textbf{Primordial State} leading to a Golden Age.

The doctrine of castes refers to the tripartite functional organization of society: the spiritual, the temporal, and the commercial. Then there is the fourth caste of the mass and beyond that the outcasts who have no fixed role in political power.

These three foundations are interrelated. The reversal of the powers of the various castes is analogous to the doctrine of cycles and result in the improper relationship between spiritual authority and political power. Any socio-political program, therefore, that claims to be based on Tradition must address in a concrete way each of these foundations.

References:
\textit{Evola Guenon De Giorgio} by Piero Di Vona

\textit{Rene Guenon et la politique} by Jean Hani

\flrightit{Posted on 2010-08-02 by Cologero}

